% concatenated from main.tex
\documentclass[a4paper, reqno]{amsart}
\usepackage[dvipsnames]{xcolor}
\usepackage{amsmath,amsthm,amssymb,relsize,mathrsfs,csquotes,microtype}
\usepackage[margin=30mm]{geometry}
%\usepackage{tikz}
% \usetikzlibrary{
%   cd,
%   calc,
%   positioning,
%   fit,
%   arrows,
%   decorations.pathreplacing,
%   decorations.markings,
%   shapes.geometric,
%   backgrounds,
%   bending
% }
%\usepackage{tikzsymbols}

\usepackage{enumitem}
\setlist{nosep,font=\normalfont,beginpenalty=10000, midpenalty=10000}

\definecolor{darkblue}{rgb}{0,0,0.6}
\newcommand{\purple}{\textcolor{purple}}
\newcommand{\red}{\textcolor{red}}
\newcommand{\blue}{\textcolor{blue}}

\usepackage[breaklinks, pdftex, ocgcolorlinks,colorlinks=true, citecolor=darkblue, filecolor=darkblue, linkcolor=darkblue, urlcolor=darkblue]{hyperref}
%\usepackage[capitalize,noabbrev]{cleveref} for some reason it does not work properly for citing definitions

% Theorems

%%%% From https://tex.stackexchange.com/questions/422/how-do-i-repeat-a-theorem-number
\makeatletter
\newtheorem*{rep@theorem}{\rep@title}
\newcommand{\newreptheorem}[2]{%
	\newenvironment{rep#1}[1]{%
		\def\rep@title{#2 \ref{##1}}%
		\begin{rep@theorem}}%
		{\end{rep@theorem}}}
\makeatother
%%%%%

\newtheorem{proposition}{Proposition}[section]
\newtheorem{theorem}[proposition]{Theorem}
\newtheorem{corollary}[proposition]{Corollary}
\newtheorem{lemma}[proposition]{Lemma}
\newreptheorem{theorem}{Theorem}
\newreptheorem{proposition}{Proposition}

% \newtheorem{thmx}{Theorem}
% \renewcommand{\thethmx}{\Alph{thmx}}

\theoremstyle{definition}
\newtheorem{definition}[proposition]{Definition}
\newtheorem{question}[proposition]{Question}
\newtheorem{problem}[proposition]{Problem}
\newtheorem{example}[proposition]{Example}
\newtheorem{conjecture}[proposition]{Conjecture}
\newtheorem{convention}[proposition]{Convention}
\newtheorem{construction}[proposition]{Construction}
%\renewtheorem{algorithm}[proposition]{Algorithm}




\theoremstyle{remark}
\newtheorem{remark}[proposition]{Remark}
\newtheorem*{claim}{Claim}
\newtheorem*{remark*}{Remark}
\newtheorem*{ackn}{Acknowledgements}

% \crefname{definition}{definition}{definitions}
% \Crefname{definition}{Definition}{Definitions}

% Definitions
\newcommand{\BDiff}{\text{BDiff}}
\newcommand{\lk}{\operatorname{lk}}
\newcommand{\sign}{\operatorname{sign}}
\newcommand{\Q}{\mathbb{Q}}
\newcommand{\N}{\mathbb{N}}
\newcommand{\T}{\mathbb{T}}
\newcommand{\CC}{\mathbb{C}}
\newcommand{\R}{\mathbb{R}}
\newcommand{\Z}{\mathbb{Z}}
\newcommand{\F}{\mathcal{F}}
\newcommand{\B}{\mathcal{B}}
\newcommand{\C}{\mathcal{C}}
\newcommand{\A}{\mathcal{A}}
\newcommand{\M}{\mathcal{M}}
\newcommand{\PD}{PD}
\newcommand{\im}{\operatorname{Im}}
\newcommand{\Id}{\operatorname{Id}}
\newcommand{\id}{\operatorname{Id}}
\newcommand{\Arf}{\operatorname{Arf}}
\newcommand{\ot}{\leftarrow}
\newcommand{\onto}{\twoheadrightarrow}
\newcommand{\into}{\hookrightarrow}
\newcommand{\ol}{\overline}
\newcommand{\ul}{\underline}
\newcommand{\wt}{\widetilde}
\newcommand{\wh}{\widehat}
\newcommand{\sm}{\setminus}
\newcommand{\Sum}{\displaystyle \sum}
\newcommand{\pri}{\mathfrak{pri}}
\newcommand{\heq}{\simeq}
\newcommand{\ks}{\operatorname{ks}}
\newcommand{\CP}{\mathbb{C}P}
\newcommand{\RP}{\mathbb{R}P}
\DeclareMathOperator{\km}{km}
\DeclareMathOperator{\Out}{Out}
\DeclareMathOperator{\Aut}{Aut}
\DeclareMathOperator{\Spin}{Spin}
\DeclareMathOperator{\GL}{GL}
\DeclareMathOperator{\pt}{pt}
\DeclareMathOperator{\pr}{pr}
\DeclareMathOperator{\TOP}{TOP}
\DeclareMathOperator{\MSTop}{MSTop}
\DeclareMathOperator{\BTOP}{BTOP}
\DeclareMathOperator{\BSTop}{BSTop}
\DeclareMathOperator{\OO}{O}
\DeclareMathOperator{\SO}{SO}
\DeclareMathOperator{\G}{G}
\DeclareMathOperator{\BG}{BG}
\DeclareMathOperator{\BO}{BO}
\DeclareMathOperator{\BSO}{BSO}
\DeclareMathOperator{\BSpin}{BSpin}
\DeclareMathOperator{\Diff}{Diff}
\DeclareMathOperator{\Hom}{Hom}
\DeclareMathOperator{\PL}{PL}
\DeclareMathOperator{\Homeo}{Homeo}
\DeclareMathOperator{\Isom}{Isom}
\DeclareMathOperator{\simp}{simp}
\DeclareMathOperator{\Tors}{Tors}
\newcommand{\Diffeo}{\operatorname{Diffeo}}

% definitions Stefan
\def\op{\operatorname}
% comments
\newcommand{\comm}[2]{%
  {\color{red}#1: #2}}

\newcommand{\commt}[2]{%
  {\color{blue}#1: #2}}

  %For splitting words

  \hyphenation{ori-en-ta-tion-pre-serv-ing} 


%defining equations
\newcommand{\defeq}{\mathrel{\mathop:}=}
\newcommand{\eqdef}{\mathrel{\mathop=}:} % Some more commands
\newcommand{\ie}{i.e.\~}
\newcommand{\Ie}{i.e.\~}
\newcommand{\eg}{e.g.\~}
\newcommand{\Eg}{e.g.\~} 
\newcommand{\Mtop}{\mathcal{M}^{\text{TOP}}}
\newcommand{\Mtopn}{\mathcal{M}^{\text{TOP}}_n}
\newcommand{\Mtopfour}{\mathcal{M}^{\text{TOP}}_4}
\newcommand{\Msm}{\mathcal{M}^{\text{Diff}}}
\newcommand{\Mpl}{\mathcal{M}^{\text{PL}}}
\newcommand{\Mtria}{\mathcal{M}^{\Delta}}
\newcommand{\Mtrian}{\mathcal{M}^{\Delta}_n}
\newcommand{\Mtriafour}{\mathcal{M}^{\Delta}_4}
\newcommand{\MM}{\mathcal{M}^{\text{sc}}_4}
\newcommand{\sumStwo}[1]{\ensuremath{\#_{#1} (S^2 \times S^2)}}

\newcommand{\ttimes}{\mathbin{%
    \ooalign{\raise1.15ex\hbox{$\scriptstyle\sim$}\cr\hidewidth$\times$\hidewidth\cr}%
    }}

\newcommand{\Stwotilde}{\ensuremath{S^2\ttimes S^2}}

\renewcommand{\~}{\nobreak\ }

\begin{document}

\title{The unsolvability of the homeomorphism problem}


\author{Stefan Friedl}
\author{Tobias Hirsch}
\author{Marc Kegel}


%Sorted addresses
\address{Fakult\"at f\"ur Mathematik, Universit\"at Regensburg, Germany}
\email{sfriedl@gmail.com}
\email{tobias.hirsch@ur.de}

\address{Universidad de Sevilla, Dpto.\ de Álgebra,
Avda.\ Reina Mercedes s/n,
41012 Sevilla, Spain}
\email{kegelmarc87@gmail.com}



\def\subjclassname{\textup{2020} Mathematics Subject Classification}
\expandafter\let\csname subjclassname@1991\endcsname=\subjclassname

\subjclass{
57-08; %Computational methods for problems pertaining to manifolds and cell complexes
03D80, %Applications of computability and recursion theory
57K40. %General topology of 4-manifolds
}
\keywords{Homeomorphism problem, algorithms, Kirby diagrams, 4-manifolds}

\date{\today} % date on first page

\begin{abstract}
In this short expository note, we give a detailed proof of Markov’s theorem on the unsolvability of the homeomorphism problem and of the existence of unrecognizable manifolds in all dimensions larger than $3$. 
\end{abstract}

\maketitle


% \comm{To Do:}{Look and cite the following sources:}

% \url{https://www.sciencedirect.com/science/chapter/bookseries/abs/pii/S0049237X08705184}

% \url{https://pagesperso.g-scop.grenoble-inp.fr/~lazarusf/Enseignement/compuTopo7.pdf}

% \comm{TH}{Both follows the argument from \cite{Markov}. The second is part of lecture notes and omits a great number of details, in particular, it does not discuss what precisely the input of the algorithm might ought to be. The first is a very detailed account, including a proper introduction into the computational aspects via Turing machines/recursive functions. It also appears to give a finite representation for smooth manifolds along the ideas MK suggested 'embed  nicely into $\R^n$ and approximate by rational points'.}


%=======================================
\section{Introduction}
\addtocontents{toc}{\protect\setcounter{tocdepth}{1}}%For hiding subsections of introduction in table of content

As a corollary of the resolution of the geometrization conjecture\~\cite{perelman1,perelman2} the homeomorphism and diffeomorphism problem is solved for closed orientable manifolds of dimension less than or equal to\~$3$, see for example\~\cite{Kuperberg}. %\footnote{More precisely, this means that there exists an algorithm that takes as input simplicial complexes whose underlying spaces are homeomorphic to connected closed orientable $n$-dimensional manifolds $M_1$ and $M_2$ for $n\leq3$ and outputs whether or not $M_1$ is homeomorphic to $M_2$. Since in these dimensions every manifold admits a unique smooth structure up to diffeomorphism\~\cite{munkres_1,munkres_2} this also decides whether or not the manifolds are diffeomorphic.}
On the contrary, a well-known theorem due to Markov\~\cite{Markov} from 1958 states that the homeomorphism problem for manifolds of dimension larger than $3$ is not solvable. For different write-ups of this result we refer to\~\cite{Stanko, Chernavsky,Kirby_Markovsthm,Gordon,Tancer,BooneHakenPoenaru} and\~\cite[Exercises\~5.1.10(c),\~5.2.2(c)]{GS}. The main idea in all these proofs and related results is to produce a solution to an undecidable problem in group theory from such a potential algorithm. 

In this expository article, we first make these undecidability statements precise and then present detailed proofs of them.
The input data of an algorithm is supposed to be a finite set. In our setting, we will work with PL-simplicial manifolds, as defined in Definition\~\ref{def:combinatorial-manifold}. (Note that these objects go by different names, at times they are also called combinatorial manifolds or stellar manifolds, see e.g.\ \cite{Glaser}.) 
Given a PL-simplicial manifold $M$, we can also consider the corresponding topological realization $|M|$ which comes with a natural PL-structure. Note that two PL-simplicial manifolds are combinatorially homeomorphic (i.e.\ stellar equivalent) if and only if the PL-manifolds given by their topological realizations are PL-homeomorphic.
%
%Recall that according to \cite[Theorems\~10.5 and\~10.6]{Munkres1966} every smooth manifold admits a smooth simplicial structure and that this smooth simplicial structure is (up to subdivision) unique up to a simplicial isomorphism. 
%Furthermore, by \cite[p.\~82]{Munkres1966} we know that any smooth simplicial structure corresponds in fact to a PL-simplicial manifold.
%Together with the Alexander--Newman theorem\~\cite{Al1930,newman} this implies that we have a well-defined map from smooth manifolds up to diffeomorphism to PL-simplicial manifolds up to PL-homeomorphisms.   ++++ SF: do we want to talk about smooth anyway? This result is due to  Munkres  (?)++++

With this background on PL-manifolds we can precisely state the unsolvability theorems.

\begin{theorem}
\label{thm:main_PL_input_and_output}
For all integers $n\geq4$ there exists no algorithm that 
\begin{itemize}
    \item takes as input connected closed orientable $n$-dimensional PL-simplicial manifolds  $M_1$ and $M_2$, and 
    \item outputs whether or not $|M_1|$ is PL-homeomorphic to $|M_2|$.
\end{itemize}
\end{theorem}

Theorem\~\ref{thm:main_PL_input_and_output} is satisfying because the in- and output are both from the category of PL-simplicial manifolds. 
The theorem is an immediate consequence of the following more general theorem,  with in- and outputs mixing different equivalent relations on manifolds.

\begin{theorem}
\label{thm:main_triangulation_input_and_arbitrary_output}
We consider one of the following five (in general different) equivalence relations on 
$n$-dimensional PL-manifolds:
\begin{enumerate}
    \item being PL-homeomorphic,
    \item being homeomorphic,
    \item being homotopy equivalent,
    \item admitting isomorphisms between all homotopy groups,
    \item admitting an isomorphism between the fundamental groups.
\end{enumerate}
For each of these equivalence relations and for all integers $n\geq4$, there exists no algorithm that 
\begin{itemize}
    \item takes as input connected closed orientable $n$-dimensional PL-simplicial manifolds $M_1$ and $M_2$% that admit a smooth structure
    , and %\comm{TH}{As I have currently written it, the argument does not provide the smooth structures. SF pointed out that this could maybe be fixed explicitly since the manifolds only have $\leq$ 2-handles and one should therefore always be able to smooth the attaching maps defining a smooth structure. 
    %
    % Abstractly, general theory of smoothing (Freedman-Quin Theorem 8.3B) also shows this as our spaces have the homotopy type of a 2-complex and PL/DIFF is 6-connected.    
    %We are unsure whether this is worthwhile to add or just confusing.}
    \item outputs whether or not $|M_1|$ is equivalent to $|M_2|$.
\end{itemize}
\end{theorem}

%Note that the restriction in the input data to manifolds that admit a smooth structure is not necessary but with this smaller input set the result is stronger. 

The set of closed orientable $n$-dimensional manifolds up to homeomorphism is countable\~\cite{Cheeger_Kister}. Similarly, it is a folk result that the set of closed orientable $n$-dimensional smooth manifolds up to diffeomorphism is countable. Thus, topological and smooth manifolds can, in principle, be used as input data for algorithms. On the other hand, we do not know of any good way to present such manifolds naturally by data that could be used as input for algorithms. Possible solutions for 4-dimensional topological manifolds are discussed in\~\cite{Freedman_Zuddas,algorithm_4_manifold}. An approach to represent smooth manifolds is discussed in\~\cite{BooneHakenPoenaru}, although it is less practical than the other approaches based on simplicial complexes. In dimensions $n\leq6$ the PL-homeomorphism class of a PL-simplicial manifold $M$ determines a smooth structure (unique up to diffeomorphism) on $|M|$, see for example \cite[Theorem 2]{Milnor_46later}.\footnote{In general there exist PL-simplicial manifolds that admit no smooth structure or several non-diffeomorphic smooth structures, see for example\~\cite{Milnor_46later}.} And thus Theorem\~\ref{thm:main_PL_input_and_output} also makes a statement about the smooth category in dimensions $n=4,5,6$. 

In dimension $4$, there exists another natural presentation method for compact smooth manifolds. Indeed any such manifold can be presented by a Kirby diagram\~\cite{GS}. From an algorithmic viewpoint, Kirby diagrams of compact 4-dimensional smooth  manifolds are equivalent to triangulations of compact 4-dimensional smooth manifolds, see  \cite[Lemma 8.2]{algorithm_4_manifold} for details. Thus Theorem\~\ref{thm:main_PL_input_and_output} implies the following corollary.

\begin{corollary}
\label{cor:main_Kirby_input_and_output_diff}
There exists no algorithm that 
\begin{itemize}
    \item takes as input Kirby diagrams\footnote{We can see a Kirby diagram as a link labelled with integers (for $2$-handles) or dots (for $1$-handles). There are several known ways to represent diagrams of links (up to planar isotopy), for example via PD codes~\cite{knotatlas,Mast}, DT codes~\cite{dt_codes}, Gau\ss\ codes~\cite{KnotInfo}, braid words~\cite{Artin}, or isosignatures~\cite{regina}.} of connected closed orientable 4-dimensional smooth manifolds $M_1$ and $M_2$, and 
    \item outputs whether or not $M_1$ is diffeomorphic to $M_2$.\qed
\end{itemize}
\end{corollary}


\subsection*{Unrecognizable manifolds}

In fact, the proof of Theorem\~\ref{thm:main_triangulation_input_and_arbitrary_output} shows something slightly stronger, namely that for no $n\geq4$ does there exist an algorithm that takes as input a connected closed orientable $n$-dimensional PL-simplicial manifold $M$ (or a Kirby diagram for $n=4$) and outputs whether or not there exists an integer $l\geq0$ such that $|M|$ is equivalent to $\#_l(S^2\times S^{n-2})$.
%there exist integers $l,m\geq0$ such that $|M|$ is equivalent to $\binoppenalty=10000 \relpenalty=10000 \#_{l-m} (S^2 \times S^{n-2}) \#_m (S^2 \ttimes S^{n-2})$.\footnote{Let $k\geq2$. The orientable $S^k$-bundles over $S^2$ are classified by $\pi_1(\operatorname{SO}(k))\cong\Z_2$ (compare Proposition 1.14 in\nobreak\ \cite{HatcherVB}), hence there exists precisely one non-trivial orientable $S^k$-bundle over $S^2$. We call this bundle $S^2\ttimes S^k$.}

Using a slightly stronger formulation of the underlying undecidability result from group theory, we prove in Section\~\ref{sec:unrecognizable} along the same lines a more general statement, saying informally that there are fixed manifolds which are unrecognizable. The precise statements are the following.

\begin{theorem}
\label{thm:pi_trivial}
For no integer $n\geq4$ does there exist an algorithm that 
\begin{itemize}
    \item takes as input a connected closed orientable $n$-dimensional PL-simplicial manifold $M$% that admits a smooth structure
    , and 
    \item outputs whether or not $\pi_1(|M|)$ is isomorphic to the trivial group.
\end{itemize}
\end{theorem}

\begin{theorem}
\label{thm:unrecognizable_different_categories}
There exists an integer $l\geq0$, such that for each of the equivalence relations $(1)$--$(3)$ from Theorem\~\ref{thm:main_triangulation_input_and_arbitrary_output} and for all integers $n\geq4$, there exists no algorithm that 
\begin{itemize}
    \item takes as input a connected closed orientable $n$-dimensional PL-simplicial manifold $M$% that admits a smooth structure
    , and 
    \item outputs whether or not $|M|$ is equivalent to $\#_l (S^2 \times S^{n-2})$.
\end{itemize}
\end{theorem}

In the same way as Corollary\~\ref{cor:main_Kirby_input_and_output_diff} is implied by Theorem\~\ref{thm:main_triangulation_input_and_arbitrary_output}, Theorem\~\ref{thm:unrecognizable_different_categories} implies the following. 

\begin{corollary}
\label{cor:unrecognizable_Kirby}
There exists an integer $l\geq0$ such that there exists no algorithm that 
\begin{itemize}
    \item takes as input a Kirby diagram of a connected closed orientable 4-dimensional smooth manifold $M$, and 
    \item outputs whether or not $M$ is diffeomorphic to $\#_l(S^2 \times S^2)$.\qed
\end{itemize}
\end{corollary}

It was shown by Novikov\~\cite{Novikov_sphere}, cf.\~\cite{Chernavsky}, that for every $n\geq5$ the $n$-sphere $S^n$ is unrecognizable. As observed in\~\cite{Chernavsky} it even follows from his construction that every connected compact $n$-dimensional PL-simplicial manifold $M$ is unrecognizable for $n\geq5$.

In recent years, there have been several attempts to determine the smallest (in terms of second homology) unrecognizable simply connected closed $4$-dimensional PL-simplicial manifold.
The strongest result to date is obtained by Tancer\~\cite{Tancer}, who proved that $\#_9(S^2\times S^2)$ is unrecognizable. 
Earlier results of Shtan'ko\~\cite{Stanko} and Gordon\~\cite{Gordon} established unrecognisability of $\#_l(S^2\times S^2)$ for $l=14$ and $l=12$, respectively; see also\~\cite{Chernavsky}. 
The recognisability of the simplest closed 4-dimensional manifold, namely the $4$-sphere $S^4$, remains unknown. In other words, the following fundamental problem is still open; see for example\~\cite{Weinberger,Kirby_Markovsthm,Gordon,Tancer}.

\begin{question}\label{conj:S4_intro}
Does there exist an algorithm that
\begin{itemize}
    \item takes as input a connected closed orientable 4-dimensional PL-simplicial manifold $M$, and
    \item decides whether or not $|M|$ is PL-homeomorphic to $S^4$?
\end{itemize}
\end{question}

\subsection*{Conventions}
With a manifold we always mean a topological manifold. If the manifold has extra attributes such as a PL-structure or a smooth structure, we mention this explicitly.

%\comm{TH}{aus den anderen paper:}
%\subsection*{A note about algorithms}
We take a pragmatic approach to the notion of algorithms. (The underlying precise notion is in terms of Turing machines~\cite{boolos_burgess_jeffrey} or equivalent models of computation.) For each algorithm, we will explain how the input is represented as finite data. We will then describe the algorithms in standard mathematical language, giving enough details that the translation to a concrete algorithmic setting is evident.

\subsection*{Acknowledgments}
The idea for this article started during the workshop on \textit{Algorithms in $4$-manifold topology}, hosted at Universit\"at Regensburg in September 2024, funded by the SFB\~1085 \emph{Higher Invariants} (Universit\"at Regensburg, funded by the DFG, ID 224262486), and organised by Stefan Friedl, Marc Kegel, and Birgit Tiefenbach. We thank the participants of the workshop for useful discussions. Special thanks go to Martin Tancer for giving a beautiful introductory talk on his work \cite{Tancer} and to Lisa Schambeck and Matthias Uschold for initial discussions about this subject.
The article was written during a week-long visit of SF and TH at the IMUS (Instituto de mathemáticas, Universidad de Sevilla) in February 2026. We thank the VII Plan Propio de Investigación y Transferencia of the University of Sevilla for funding of the visit and the IMUS for providing office space.

\subsection*{Individual grant support}
MK is supported by a Ram\'on y Cajal grant \mbox{(RYC2023-043251-I)} and PID2024-157173NB-I00 funded by MCIN/AEI/10.13039/501100011033, by ESF+, and by FEDER, EU; and by a VII Plan Propio de Investigación y Transferencia (SOL2025-36103) of the University of Sevilla.


% \section{Constructing manifolds with given fundamental group}
% We reduce the above decision problems of manifolds to the triviality problem for groups using the following construction.

% \begin{definition}
% \label{def:construction}
% Let $P=\langle g_1,\dots,g_k \mid r_1,\dots,r_l\rangle$ be a finite presentation of a group and $n\geq4$. Let $\mathcal W_P^n$ be the set of diffeomorphism classes of smooth $(n+1)$-manifolds that can be obtained in the following way.
% \begin{enumerate}
% \item \label{def:MP-constr-2} First build a smooth manifold $W_1$ by attaching $k$ many 1-handles via orientation-preserving attaching maps to $D^{n+1}$. 

% \noindent Note that there exists a canonical isomorphism between $\pi_1(W_1)$ and $\langle g_1,\dots,g_k\rangle$ and that the inclusion $\partial W_1\to W_1$ induces an isomorphism on fundamental groups because we only attach handles of codimension at least three. 
% \item \label{def:MP-constr-3} 
% Attach $l$ many 2-handles to $W_1$ along attaching maps $\varphi_1,\dots,\varphi_l\colon\partial D^2\times D^{n-2}\to \partial W_1$ such that $\varphi_i(\partial D^2\times\{0\})$ represents $r_i\in \langle g_1,\dots,g_k\rangle$ under the above isomorphism. 
% \item \label{def:MP-constr-4} Finally take the $k$-fold boundary connected sum of $W_2$ with $S^2\times D^{n-1}$.
% \end{enumerate}
% \end{definition}

% \begin{lemma}\label{lem:welldefined}
%     Let $P$ be a finite group presentation with $l$ relations and $n\geq4$.
%     \begin{enumerate}
%         \item The set $\mathcal W_P^n$ contains at most $2^l$ elements.     
%         %The set $\mathcal W_P^n$ is well-defined, i.e.\ does not depend on the choices made in the construction. It has at most $2^l$ elements where $l$ is the number of relators of $P$. Marc: I changed the statement. Since it is just defined as diffeo classes it is well-defined automatically.
%         \item For each $W\in\mathcal W_P^n$, the inclusion $\partial W\to W$ induces an isomorphism on fundamental groups and $P$ is a presentation for these groups. 
%     \end{enumerate}
% \end{lemma}
% \begin{proof}
% (1) In the construction of $W\in\mathcal W_P^n$ choices up to diffeomorphism are only made in Step \ref{def:MP-constr-3}. Here, for each $i\in\{1,\dots,l\}$ a choice of a $\varphi_i$ corresponds to 
%         \begin{itemize}
%             \item a smooth submanifold $\gamma_i=\varphi_i(\partial D^2\times\{0\})\subseteq\partial W_1$ diffeomorphic to $S^1$ representing $r_i\in \langle g_1,\dots,g_k\rangle$ under the above isomorphism, and
%             \item a trivialization of a normal bundle of $\gamma_i$.
%         \end{itemize}
%         By the uniqueness of tubular neighbourhoods  (see for example \cite{friedlmonster}) there are two choices for $\varphi_i$ up to smooth isotopy since
%         \begin{itemize}
%             \item by general position in dimension $n\geq4$ homotopy of curves implies smooth isotopy, so $\gamma_i$ is determined up to smooth isotopy, and
%             \item the normal bundle of $\gamma_i\subseteq\partial W_1$ is a rank $(n-1)$ real vector bundle, implying that its trivializations are classified by $\pi_1(\operatorname{SO}(n-1))\cong\Z_2$. 
%         \end{itemize}
%         Since isotopic handle attachments lead to diffeomorphic manifolds, the claim follows. 
   
%         (2) By construction $\pi_1(W)$ is presented by $P$. The inclusion $\partial W\to W$ induces an isomorphism on fundamental groups since $W$ has a handle decomposition where all handles have codimension at least 3. Note that the boundary connected sum with $S^2\times D^{n-1}$ can be realized by attaching a $2$-handle along a homotopically trivial curve\~\cite{GS}. 
% \end{proof}

% The following lemma is the main idea in Markov's proof\~\cite{Markov} and says that the diffeomorphism type of a manifold in $\mathcal W_P$ is detected by its fundamental group.

% \begin{lemma}\label{lem:P}
%     % Let $n\geq4$ and consider one of the following equivalence relations on smooth $n$-dimensional manifolds:
%     % \begin{enumerate}\setcounter{enumi}{-1}
%     %     \item being diffeomorphic,
%     %     \item being PL-homeomorphic, \comm{TH}{check it has been explained how to interprete this}
%     %     \item being homeomorphic,
%     %     \item being homotopy equivalent,
%     %     \item admitting isomorphisms between all homotopy groups,
%     %     \item admitting an isomorphism between the fundamental groups.
%     % \end{enumerate}
%     %Let $P$ be a finite group presentation and $W\in\mathcal W_P^n$. Then $P$ represents the trivial group if and only if 
%     Let $n\geq4$. If $P$ is a presentation of the trivial group with $l$ relations, then for every $W\in\mathcal W^n_P$, there exists an $m\in\N_0$ such that $\partial W$ is diffeomorphic to 
%     \[
%     \#_{l-m} (S^2 \times S^{n-2})\#_m(S^2\ttimes S^{n-2}).
%     \] 
% \end{lemma}

% \begin{proof}
%     %As the equivalence relations imply each other in the given order, it suffices to prove that $P$ representing the trivial group implies the statement when \enquote{equivalent} means \enquote{diffeomorphic}:
%     %The \enquote{only if}-direction is clear, so it remains to consider the \enquote{if}-direction:
%     By Lemma \ref{lem:welldefined} $\pi_1(W)\cong\pi_1(\partial W)$ are trivial. Step \ref{def:MP-constr-4} in the construction of $W$ corresponds to attaching $k$ many 2-handles. Hence, $\pi_1(\partial W_2)\cong\pi_1(\partial W)$ is trivial, implying that the attaching sphere of such a $2$-handle is homotopic, hence isotopic, to a loop that cancels one of the $1$-handles 
%     attached in Step \ref{def:MP-constr-2}. 
%     Thus, $W$ is diffeomorphic to $D^{n+1}$ with $l$ 
%     many $2$-handles attached. As discussed in Lemma \ref{lem:welldefined} there are precisely two choices for each 2-handle up to isotopy, which in the boundary lead to a summand of $S^2 \times S^{n-2}$ or $S^2\ttimes S^{n-2}$.
% \end{proof}

% The following lemma is not needed for the proof of Theorem\~\ref{thm:main_triangulation_input_and_arbitrary_output} but is needed for the strengthenings, Theorem\~\ref{thm:pi_trivial} and Theorem\~\ref{thm:unrecognizable_different_categories}.

% \begin{lemma}\label{lem:P_strong}
% Let $n\geq4$. If $P$ is a presentation of the trivial group with $l$ relations, then there exists a manifold $W\in\mathcal W^n_P$ such that $\partial W$ is diffeomorphic to 
%     \[
%     \#_{l} (S^2 \times S^{n-2}).
%     \] 
% \end{lemma}

% For the proof of Lemma\~\ref{lem:P_strong} we need the following two results.

% \begin{lemma}\label{lem:embedding}
%     Let $P$ be a finite group presentation. Then there exists at least one $W\in\mathcal W_P^n$ that embeds in $\R^{n+1}$. In particular, $\partial W$ has stably trivial tangent bundle and thus vanishing second Stiefel--Whitney class.
%     %with trivial tangent bundle / trivial seconds Stiefel--Whitney class.
% \end{lemma}

% \begin{proof}
% %It suffices to show that in the above construction at least one choice leads to a codimension-zero submanifold of  $\R^5$.

% %We attach the 1-handles in the standard way to obtain a codimension-zero submanifold $W_1$ of $\R^5$. Let $X_1$ be  the complement of the interior of $W_1$ viewed as a submanifold of $S^5=\R^5\cup \{\infty\}$. Note that $X_1$ admits a handle decomposition with one 0-handle and $k$ 3-handles. In particular, we see that $X_1$, and thus also $X_1\cap \R^5$ is simply connected.
% For every $1$-handle $h^i_1$ of $W_1$ we can attach a $2$-handle $h_2^i$ to $W_1$ such that the attaching sphere of the $2$-handle $h_2^i$ intersects the belt sphere of the $1$-handle $h_1^i$ transversely in a single point. Thus the handles cancel and we get a handle decomposition of $D^5$ to which we can glue an $(n+1)$-handle to obtain $S^{n+1}$. This describes an embedding of $W_1$ into $S^{n+1}$. By considering the dual handle decomposition, the exterior of $W_1$ in $S^{n+1}$ can be obtained by attaching $k$ $3$-handles to a $0$-handle and thus is simply connected. By seeing $S^{n+1}$ as $\R^{n+1}\cup \{\infty\}$, we also get an embedding of $W_1$ into $\R^{n+1}$ with simply connected exterior $X_1$.

% Now let $\gamma_1,\dots,\gamma_l$ in $\partial W_1$ be smoothly embedded curves that correspond to $r_1,\dots,r_l$. Since the exterior
% $X_1$ of $W_1$ in $\R^{n+1}$ is simply connected and since $\dim(X_1)=n+1>2\cdot 2$, there exist proper smooth embeddings $\varphi_1,\dots,\varphi_l\colon D^2 \to X_1\subseteq \R^{n+1} $
% such that for each $i$ we have $\varphi_i(S^1)=\gamma_i$.
% We choose tubular neighbourhoods of the proper submanifolds $\varphi_i(D^2)\subseteq X_1$, $i=1,\dots,l$. Since $D^2$ and $\R^{n+1}$ are orientable, we see that the tubular neighbourhoods are trivial. We now use these trivial tubular neighbourhoods to attach the $l$ 2-handles forming $W_2$, which corresponds to one of the $2^l$ choices and is now by construction a codimension-zero submanifold of\~$\R^{n+1}$ with exterior\~$X_2$.

% Finally, we choose $k$ disjoint homotopically trivial curves $\alpha_1,\dots,\alpha_k$ in $\partial W_2$ and proper smooth embeddings $\psi_1,\dots,\psi_k\colon D^2 \to X_2 \subseteq \R^{n+1} $
% such that for each $i$ we have $\psi_i(S^1)=\alpha_i$. As proper submanifolds of $X_2\subset \R^{n+1}$ the $\varphi_i(D^2)$ have trivial tubular neighbourhood which we can use to attach $k$ $2$-handles along the $\alpha_i$. Since the tubular neighbourhoods are trivial, these $2$-handle attachments correspond to the $k$-fold boundary connected sum with $S^2\times D^{n-1}$. Thus we have realized $W$ as a codimension-zero submanifold of\~$\R^{n+1}$.
% \end{proof}

% The following lemma is well-known. We provide a proof for completeness.

% \begin{lemma}\label{lem:w2}
%     Let $k\geq2$. The second Stiefel--Whitney class of the manifold $S^2\ttimes S^k$ is non-vanishing. 
% \end{lemma}

% \begin{proof}
%     % We proceed through several intermediate claims.
%     % \begin{claim}
%     %     There exists precisely one non-trivial $\R^{k+1}$-bundles over $S^2$ and its unit sphere bundle is $S^2\ttimes S^k$.
%     % \end{claim}
%     % \begin{proof}\renewcommand{\qedsymbol}{\ensuremath\boxplus}
%     %     The $\R^{k+1}$-bundles over $S^2$ are classified by $\pi_2(\operatorname B\operatorname{SO}(k))$ where $\operatorname B\operatorname{SO}(k)$ is the classifying space of the topological group $\operatorname{SO}(k)$. The connecting homomorphism $\pi_2(\operatorname B\operatorname{SO}(k))\to\pi_1(\operatorname{SO}(k))$ in the long exact sequence of the fibration $G\to\operatorname E\operatorname{SO}(k)\to\operatorname B\operatorname{SO}(k)$ is an isomorphism. On the level of bundles it corresponds to taking the unit sphere bundle. 
%     % \end{proof}
%     Let $p\colon S^2\ttimes\R^{k+1}\to S^2$ be the associated vector bundle to the sphere bundle $S^2\ttimes S^k\to S^2$. This bundle is classified by an element $f\in\pi_2(\operatorname B\operatorname{SO}(k))$ where $\operatorname B\operatorname{SO}(k)$ is the classifying space of the topological group $\operatorname{SO}(k)$. We have a sequence of natural isomorphisms
%     \[
%     \Phi\colon\pi_2(\operatorname B\operatorname{SO}(k))\to\operatorname H_2(\operatorname B\operatorname{SO}(k);\Z_2)\to\operatorname H_2(\operatorname B\operatorname{SO}(k);\Z_2)^{**}\to\operatorname H^2(\operatorname B\operatorname{SO}(k);\Z_2)^*
%     \]
%     where
%     \begin{itemize}
%         \item the first is given by the Hurewicz theorem (note that $\pi_1(\operatorname B\operatorname{SO}(k))\cong\pi_0(\operatorname{SO}(k))\cong1$),
%         \item the second is the double-dual space isomorphism over $\Z_2$,
%         \item and the third is the dual of the evaluation isomorphism from the universal coefficient theorem over the field $\Z_2$.
%     \end{itemize}
%     Then $\Phi(f)(\omega_2)=\langle\omega_2,f_*[S^2]\rangle=\langle f^*\omega_2,[S^2]\rangle$ where $\omega_2\in\operatorname H^2(\operatorname B\operatorname{SO}(k);\Z_2)$ is the generator and $f^*\omega_2=w_2(p)\in\operatorname H^2(S^2)$ (compare Proposition 3.12 in \cite{HatcherVB}). Since $\Phi$ is an isomorphism and $f$ is non-trivial it follows that $w_2(p)$ is non-vanishing.

%     Next we calculate the Stiefel--Whitney class $w_2(S^2\ttimes\R^{k+1})$ of the total space of $p$ as a manifold in its own right. Its tangent space is given by
%     \[
%     T(S^2\ttimes\R^{k+1})\cong p^*(TS^2)\oplus p^*(S^2\ttimes\R^{k+1}).
%     \]
%     Hence, $w_2(S^2\ttimes\R^{k+1})=p^*w_2(p)$. Since $p^*\colon\operatorname H^2(S^2;\Z_2)\to\operatorname H^2(S^2\ttimes\R^{k+1};\Z_2)$ is injective by the Gysin-sequence, $w_2(S^2\ttimes\R^{k+1})$ is non-vanishing.    
%     The inclusion $S^2\ttimes S^k\to S^2\ttimes\R^{k+1}$ is a homotopy equivalence, so the claim follows. 
% \end{proof}

% \begin{proof}[Proof of Lemma \ref{lem:P_strong}]
% By Lemma\~\ref{lem:embedding} there exist a manifold $W$ in $\mathcal{W}_P^n$ that embeds as a codimension-zero submanifold into $\R^{n+1}$. In particular, the tangent bundle of $\partial W$ is stably trivial and thus the second Stiefel--Whitney class of $\partial W$ is vanishing. 

% On the other hand, by Lemma\~\ref{lem:P} there exist an integer $m\geq0$ such that $\partial W$ is diffeomorphic to $\#_{l-m} (S^2 \times S^{n-2})\#_m(S^2\ttimes S^{n-2})$. Since the second Stiefel--Whitney class of $S^2\times S^{n-2}$ is trivial and the second Stiefel--Whitney class is additive under connected sum, it follows from Lemma\~\ref{lem:w2} that $m=0$ as the second Stiefel--Whitney class of $\partial W$ vanishes.
% \end{proof}


%===============================================
% \section{From presentations to PL-simplicial manifolds}\label{section:presentation-to-pl}
% In this section, we will assume that the reader has some familiarity with concepts from combinatorial topology and PL-topology. We refer to \cite{zeeman,hudson,Glaser,rourkesanderson,friedlmonster} for details.  



% \begin{definition}\label{def:combinatorial-manifold}
% \begin{enumerate}
% \item We say that two finite abstract simplicial complexes $K$ and $L$ are \emph{combinatorially homeomorphic} if $K$ and $L$ are stellar equivalent, i.e.\ if they are related by  a sequence of stellar subdivisions, stellar welds and simplicial isomorphisms. 
% \item Let $n\in \N_0$. 
%  We say that a finite abstract simplicial complex is a 
% \emph{combinatorial $n$-ball} if it is combinatorially homeomorphic to the abstract simplicial complex $\Delta^n\defeq\{\{0,\dots,n\},\mathcal{P}(\{0,\dots,n\})\}$
% and we say that a finite abstract simplicial complex is a 
% \emph{combinatorial $n$-sphere} if it is combinatorially homeomorphic to the abstract simplicial complex $\partial \Delta^n\defeq\{\{0,\dots,n\},\mathcal{P}(\{0,\dots,n\})\setminus \{0,\dots,n\}\}$. 
% \item Let $K$ be a finite abstract simplicial complex. 
% We say that $K$ is an \emph{$n$-dimensional \mbox{PL-simplicial} manifold} if for every vertex $v$ the link $\op{Lk}(K,v)$ is a combinatorial $(n-1)$-sphere or a combinatorial $(n-1)$-ball.
% \item The boundary $\partial K$ of an $n$-di\-men\-sional PL-simplicial manifold $K$ is defined as the union of all simplices $s$ such that  the link $\op{Lk}(K,s)$ is a combinatorial ball. Note that  by \cite[Corollary\~II.4]{Glaser} we know that  $\partial K$ is again a PL-simplicial manifold.
% We say $K$ is closed if $K$ is finite and $\partial K=\emptyset$.
% \end{enumerate}
% \end{definition}

% Note that it follows from the Alexander--Newman theorem\~\cite[Theorem\~1.5]{Al1930}, \cite{newman}, 
% that the natural map from PL-simplicial manifolds to PL-manifolds induces a bijection of the stellar equivalence classes of PL-simplicial manifolds and PL-homeomorphism classes of PL-manifolds.



% In Definition\~\ref{def:construction} we associated to each group presentation $P=\langle g_1,\dots,g_k \mid r_1,\dots,r_l\rangle$  a set $\mathcal{W}^n_P$ of $(n+1)$-dimensional \emph{smooth} manifolds whose fundamental groups are presented by $P$. In the following we will sketch how to algorithmically construct a PL-simplicial manifold realizing one element in $\mathcal{W}_P^n$. We do not claim that our construction gives a full algorithm, but we hope that it convinces the reader that one could in principle make this completely algorithmic.
% In the construction we make use of the following observation: If given two finite abstract simplicial complexes we know for theoretical reasons that a combinatorial isomorphism exists, then one can also find it by going through all (countably many! \comm{TH}{countability does not suffice, as there do exist countable sets which are not the iterative output of an algorithm (eg. the set of all Turing machines that do not halt). I am unsure how do phrase this, perhaps we should just drop the bracket entirely but this feels like cheating.}) possibilities. 

% \begin{lemma}\label{lem:pl-for-wp}
% For any $n\geq4$, there exist an algorithm that
% \begin{itemize}
%     \item takes as input a finite group presentation $P$, and
%     \item outputs a finite PL-simplicial manifold whose topological realization is PL-homeomorphic to one element in $\mathcal{W}_P^n$.
% \end{itemize}
% \end{lemma}

% \begin{proof}
% The idea of the proof is to perform Definition\~\ref{def:construction} with PL-simplicial manifolds and then follow the logic of 
% Lemma\~\ref{lem:welldefined}. The literature that relates smooth constructions to PL-constructions is unfortunately not very extensive. So we will not attempt to provide full details. 

% Let $P=\langle g_1,\dots,g_k \mid r_1,\dots,r_l\rangle$ be a finite presentation of a group. We build an $(n+1)$-dimensional PL-simplicial manifold
% $X_P$ as follows.
% \begin{enumerate}
% \item   We attach $k$ many 1-handles via orientation-preserving simplicial maps 
% to a barycentric subdivision of the simplicial ball
% $\Delta^{n+1}=\{\{0,\dots,n+1\},\mathcal{P}(\{0,\dots,n+1\})\}$. 
% Let $X_1$ be the resulting PL-simplicial manifold. 
% Note that we have a canonical identification $\pi_1(X_1)\cong\langle g_1,\dots,g_k\rangle$.

% \item Possibly after a further barycentric subdivision we can pick disjoint closed oriented simplicial  loops $C_1,\dots,C_l\subset \partial X_1$ 
% such that 
% $C_1,\dots,C_l$ represent the free homotopy classes 
% of $r_1,\dots,r_l\in \pi_1(X_1)\cong \langle g_1,\dots,g_k\rangle$.

% Next we can consider the regular neighbourhoods $N_1,\dots,N_l$
% of $C_1,\dots,C_l\subset \partial X_1$ as in \cite[p.\~33]{rourkesanderson} or \cite[Chapter\~III, p.\~22]{zeeman}. Recall that these can be constructed by taking the double barycentric subdivision of $X_1$ and by considering the union of the stars of all simplices in $N_1,\dots,N_l$.

% It follows from \cite[Theorems\~3.8 and 3.11]{rourkesanderson} \comm{TH+SF}{we agree that this is not really true} that each $N_i$ is combinatorially homeomorphic
% to the PL-simplicial manifold $\partial \Delta^2\times \Delta^{n-1}$. It follows from \cite[Lemma\~II.3]{Glaser} (applied twice, first to $\Delta^2\times \partial \Delta^{n-1}$ and then to $\Delta^2\times \Delta^n$)
% that for each $i\in \{1,\dots,n\}$ there exists a PL-simplicial manifold $Y_i$ that is 
% combinatorially homeomorphic to $\Delta^2\times \Delta^{n-1}$
% and such that the combinatorial structure 
% on $\partial \Delta^2\times \Delta^{n-1}$ is isomorphic to $N_i$.
% We can now use this isomorphism to attach combinatorial 2-handles to $X_1$ along $N_1,\dots,N_l$. We call the resulting combinatorial  manifold $X_2$.
% \item Finally we  perform the combinatorial boundary connected sum of $X_2$ with $k$ copies of $\partial \Delta^3\times \Delta ^{n-2}$ to obtain a PL-simplicial manifold $X_P$. Here the boundary connected sum is defined as follows:
% We perform two more barycentric subdivisions and then we glue along an $n$-simplex along the boundary.
% \end{enumerate}
% By similar arguments as in Lemma\~\ref{lem:welldefined} it follows that the constructed PL-simplicial manifold $X_P$ whose topological realization is PL-homeomorphic to one manifold from $\mathcal{W}_P^n$.
% \end{proof}

% The following lemma is a strengthening of Lemma\~\ref{lem:pl-for-wp} and allows us to algorithmically construct PL-simplicial manifolds for all elements in $\mathcal{W}_P^n$. For the proof of Theorem\~\ref{thm:main_triangulation_input_and_arbitrary_output} we only need Lemma\~\ref{lem:pl-for-wp}, but in the proof of Theorem\~\ref{thm:pi_trivial} and Theorem\~\ref{thm:unrecognizable_different_categories} we will also use Lemma\~\ref{lem:all}.

% \begin{lemma}\label{lem:all}
% For any $n\geq4$, there exist an algorithm that
% \begin{itemize}
%     \item takes as input a  finite group presentation $P$, and
%     \item outputs a finite list of  finite PL-simplicial manifolds $\mathcal{X}_P^n$, such that each manifold in $\mathcal{W}_P^{n}$ %\comm{TH}{should this really be $n+1$ and $n\geq4$?Marc: I have corrected it.} 
%     is PL-homeomorphic to the topological realization of at least one element in $\mathcal{X}_P^n$.
% \end{itemize} 
% \end{lemma}

% For the proof of Lemma\~\ref{lem:all} we need the following result.

% \begin{lemma}\label{lem:realizing_diffeo}
% Let $M$ be a compact smooth manifold that is equipped with a smooth simplicial structure $S$ and let $f\colon M\to M$  be a diffeomorphism.
% Then there exist subdivisions $S'$ and $S''$ of $S$ 
% and a simplicial isomorphism $\varphi\colon S'\to S''$
% such that $|\varphi|$ is homeotopic to $f$.    
% \end{lemma}

% \begin{proof}
% This lemma follows from the uniqueness of smooth simplicial structures on smooth manifolds (see \cite[Theorem\~10.5]{Munkres1966}) applied to the smooth simplicial structures $S$ and $f^*S$. 
% \end{proof}

% In our application we will apply the lemma to $M=D^n\times S^1$, 
% a fixed smooth simplicial structure on $D^n\times S^1$ and the diffeomorphism $D^n\to D^n\times S^1$ that is given by twisting along $S^1$. 

% \begin{proof}[Proof of Lemma\~\ref{lem:all}]
%     \comm{To do}{for Stefan}
% \end{proof}


%new idea==============
\section{From group presentations to PL-simplicial manifolds}\label{section:presentation-to-pl}
In this section, we will assume that the reader has some familiarity with concepts from combinatorial topology and PL-topology. We refer to \cite{zeeman,hudson,Glaser,rourkesandersonnew,friedlmonster} for details.  



\begin{definition}\label{def:combinatorial-manifold}\hfill
\begin{enumerate}
\item Two finite abstract simplicial complexes $K$ and $L$ are \emph{combinatorially homeomorphic} if $K$ and $L$ are stellar equivalent, i.e.\ if they are related by  a sequence of stellar subdivisions, stellar welds and simplicial isomorphisms. 
\item Let $n\in \N_0$. 
 A finite abstract simplicial complex is a 
\emph{combinatorial $n$-ball} if it is combinatorially homeomorphic to the abstract simplicial complex \mbox{$\Delta^n\defeq\{\{0,\dots,n\},\mathcal{P}(\{0,\dots,n\})\}$} and a 
\emph{combinatorial $n$-sphere} if it is combinatorially homeomorphic to the abstract simplicial complex $\partial \Delta^n\defeq\{\{0,\dots,n\},\mathcal{P}(\{0,\dots,n\})\setminus \{0,\dots,n\}\}$. 
\item A abstract simplicial complex $K$ is an \emph{$n$-dimensional \mbox{PL-simplicial} manifold} if for every vertex $v$ the link $\op{Lk}(K,v)$ is a combinatorial $(n-1)$-sphere or a combinatorial $(n-1)$-ball.
\item The boundary $\partial K$ of an $n$-di\-men\-sional PL-simplicial manifold $K$ is the union of all simplices $s$ such that  the link $\op{Lk}(K,s)$ is a combinatorial ball. By \cite[Corollary\~II.4]{Glaser} we know that $\partial K$ is again a PL-simplicial manifold.
We say $K$ is \emph{closed} if $K$ is finite and $\partial K=\emptyset$.
\item An \emph{$n$-dimensional PL-manifold} (note the missing \enquote{simplicial}) is a topological space together with a homeomorphism to the topological realization of an $n$-dimensional PL-simplicial manifold.
\end{enumerate}
\end{definition}




In \cite[Chapter 6]{rourkesandersonnew} a theory of PL-handle decompositions is developed. When adapted to our setting, it essentially parallels the well-known theory of handle decompositions for smooth manifolds, i.e.\ a $k$-handle on a PL-simplicial manifold $W$ is a $\Delta^k\times\Delta^{n-k}$ attached along a \mbox{PL-embedding} $\partial\Delta^k\times\Delta^{n-k}\to\partial W$. For technical reasons we choose the PL-structure on $\Delta^k\times\Delta^{n-k}$ that is obtained by taking the product PL-structure after barycentric subdivision of $\Delta^k$. This will be needed to ensure that the $k$-handle collapses onto its core. Also note that \cite{rourkesandersonnew} deals with PL-manifolds PL-embedded in some $\R^n$. While we in principle work with abstract manifolds, it will later be necessary to embed them. The treatment in \cite{rourkesandersonnew} then ensures that the necessary handle slides and cancellations are still possible.\pagebreak

Using handle attachments we associate PL-simplicial manifolds to a finite group presentation.  
\begin{definition}\label{dfn:relalizeP}
    Let  $P=\langle g_1,\dots, g_k\mid r_1,\dots, r_l\rangle$ be a finite group presentation and $n\geq4$. A handle decomposition on an $(n+1)$-dimensional PL-simplicial manifold $W$ \emph{realizes $P$} if it is obtained in the following manner:
    \begin{itemize}
        \item First build a PL-simplicial manifold $W_1$ by attaching $k$ many PL 1-handles via orientation-preserving attaching maps to a PL-ball $D^{n+1}$. 

        Note that up to reordering there exists a preferred isomorphism between $\pi_1(W_1)$ and $\langle g_1,\dots, g_k\rangle$ and that the inclusion $\partial W_1\to W_1$ induces an isomorphism on fundamental groups as we only attach handles of codimension at least 3. 
        \item Now attach $l$ many PL 2-handles to $W_1$ via attaching maps $\phi_1,\dots,\phi_l\colon\partial\Delta^2\times \Delta^{n-1}\to\partial W_1$ where $\phi_i(\partial \Delta^2\times\{0\})$ represents $r_i\in\langle g_1,\dots,g_k\rangle$ under the above isomorphism to form $W_2$
        \item Finally attach $k$ many additional 2-handles to $W_2$. % via attaching maps $\psi_1,\dots,\psi_k\colon\partial\Delta^2\times \Delta^{n-1}\to\partial W_1$such that $\psi_i(\partial \Delta^2\times\{0\})$ is null-homotopic in $\partial W_2$ ++++ delete +++.
        %\item Finally take the $k$-fold boundary connected sum of $W_2$ with $S^2\times D^{n-1}$.    \comm{TH}{I dont see where we use in the current argument that these handles are attached trivially. We use them for cancelling the 1-handles and there I dont think this matters.}
    \end{itemize}
\end{definition}
% \begin{lemma}\label{lem:pi1P}
%     Let $P$ be a finite presentation, $n\geq4$ and $W$ and $(n+1)$-dimensional PL-simplicial manifold that allows a handle decomposition realizing $P$. Then $\pi_1(W)\cong\pi_(\p W)$ is presented by $P$.
% \end{lemma}
% \begin{proof}
%     By construction, $\pi_1(W)$ is presented by $P$. The inclusion $\partial W\to W$ induces an isomorphism on fundamental groups since $W$ has a handle decomposition where all handles have codimension at least 3. 
% \end{proof}

To apply this construction in our argument, it is important that we can find such a manifold realizing a given group presentation algorithmically. To do this we recall the following terminology:
\begin{definition}
    A set $X$ is \emph{recursively enumerable} if there exists an algorithm that takes as input a natural number and outputs an element of $X$ such that every element of $X$ is the output of at least one natural number.
\end{definition}
It is part of this definition that the elements of $S$ are of a form that can be the output of an algorithm. Observe the following algorithmic lemma using a technique known as \enquote{dovetailing}, compare \cite[Corollary 7.15]{boolos_burgess_jeffrey}.
\begin{lemma}\label{lem:recenum}
    Let $X$ be a recursively enumerable set and $\mathcal A$ be an algorithm that takes as input an element $x\in X$. The subset $Y\subseteq X$ of elements on which $\mathcal A$ halts is recursively enumerable. 
\end{lemma}
\begin{proof}
    The algorithm $\mathcal A$ consists of a sequence of computational steps. Construct a recursive enumeration of $Y$ as follows:
    \begin{itemize}
        \item Run step 1 of $\mathcal A$ on the 1\textsuperscript{st} output of the recursive enumeration of $X$.
        \item Run step 1 of $\mathcal A$ on the 2\textsuperscript{nd} output of the recursive enumeration of $X$.
        \item Run step 2 of $\mathcal A$ on the 1\textsuperscript{st} output of the recursive enumeration of $X$.
        \item Run step 1 of $\mathcal A$ on the 3\textsuperscript{rd} output of the recursive enumeration of $X$.
        \item Run step 2 of $\mathcal A$ on the 2\textsuperscript{nd} output of the recursive enumeration of $X$.
        \item Run step 3 of $\mathcal A$ on the 1\textsuperscript{st} output of the recursive enumeration of $X$.
        %\item Run step 1 of $\mathcal A$ on the 4\textsuperscript{th} output of the recursive enumeration of $X$.
        \item $\dots$
    \end{itemize}
    If this step-wise execution of $\mathcal A$ halts on an output of the recursive enumeration of $X$, output it. 
\end{proof}
We can algorithmically find a manifold realizing $P$ that embeds into $\R^{n+1}$ by first enumerating a larger set and then narrowing down our search using the previous lemma.
\begin{lemma}\label{lem:algrelP}
    For every $n\geq4$, there exists an algorithm that
    \begin{itemize}
    \item takes as input a  finite group presentation $P$, and
    \item outputs an $(n+1)$-dimensional PL-simplicial manifold $W$ which PL-embeds into $\mathbb R^{n+1}$ and allows a handle decomposition realizing $P$.
\end{itemize} 
\end{lemma}
\begin{proof}
    A PL-simplicial manifold $W$ PL-embeds into $\R^{n+1}$ if and only if there exists a PL-simplicial manifold PL-homeomorphic to $S^{n+1}$ with $W$ as a subcomplex (see \mbox{\cite[Theorem\~2.14]{rourkesandersonnew}}). The set of compact PL-simplicial manifolds (up to simplicial isomorphism) that are PL-hom\-e\-o\-mor\-phic to $S^{n+1}$ is recursively enumerable. This is obtained by starting with some $S^{n+1}$ and systematically applying all (finitely many) possible stellar subdivisions and stellar welds, then iterating this process on the results. As each PL-structure on $S^{n+1}$ is finite, each such PL-simplicial manifold has finitely many subcomplexes. It follows that the set $X$ of subcomplexes (up to simplicial isomorphism) of PL-simplicial manifolds PL-homeomorphic to $S^{n+1}$ is recursively enumerable.   
    
    We now apply Lemma\~\ref{lem:recenum} with the following algorithms:
    \begin{itemize}
        \item There exists an algorithm with input a finite abstract simplicial complex $W$ that halts if and only if $W$ is a compact $(n+1)$-dimensional PL-simplicial manifold. This again follows by systematically applying stellar subdivisions and stellar welds to the links of the finitely many vertices of $W$ to check if they are combinatorial $n$-spheres or combinatorial $n$-balls. If this is the case, the process will eventually stop.
        \pagebreak
        \item There exists an algorithm with input a compact PL-simplicial manifold $W$ that halts if and only if $W$ allows a handle decomposition realizing $P$. This is the case as there are only finitely many possibilities for the location of the handles. 
    \end{itemize}
    Thus the set of $(n+1)$-dimensional PL-simplicial manifolds (up to simplicial isomorphism) that PL-embed into $\R^{n+1}$ and allow a handle decomposition realizing $P$ is recursively enumerable.
    
    It remains to show that it is non-empty: We realize the presentation complex $X$ of $P$ as a 2-dimensional finite abstract simplicial complex. It therefore PL-embeds into $\R^{n+1}$ for $n\geq2\cdot2$ by a general position argument \cite[Theorem\~5.4]{rourkesandersonnew}. Hence, there exists a PL-simplicial structure on $\R^{n+1}$ such that $X$ arises as a subcomplex (see \cite[Theorem\~2.14]{rourkesandersonnew}). A regular neighbourhood of $X$ in it has a handle decomposition realizing $P$ by \cite[Proposition\~6.9]{rourkesandersonnew}.
    % By systematically going through all possibilities, we can construct the following algorithms:
    % \begin{enumerate}
    %     \item An algorithm with input a natural number and output a compact $(n+1)$-dimensional PL-simplicial manifold such that every PL-simplicial manifold is the output of at least one natural number.        
    %     \item An algorithm with input a compact PL-simplicial manifold $W$ that halts if and only if $W$ PL-embeds into $\R^{n+1}$ with vertices in $\Q^{n+1}$.
    %     \item An algorithm with input a compact PL-simplicial manifold $W$ that halts if and only if $W$ allows a handle decomposition realizing $P$.
    % \end{enumerate}
    % Now construct a new algorithm dovetailing these algorithms by alternating them in the manner of the usual enumeration of $\N^2$: 
    % \begin{itemize}
    % \item Run step 1 of (2) and (3) on the 1st output of (1),
    % \item Run step 1 of (2) and (3) on the 2nd output of (1),
    % \item Run step 2 of (2) and (3) on the 1st output of (1), 
    % \item Run step 1 of (2) and (3) on the 3rd output of (1), 
    % \item Run step 2 of (2) and (3) on the 2nd output of (1), 
    % \item Run step 3 of (2) and (3) on the 1st output of (1), 
    % \item Run step 1 of (2) and (3) on the 4th output of (1), 
    % \item $\dots$
    % \end{itemize} The resulting algorithm systematically checks all possibilities for a PL-simplicial manifold $W$ which PL-embeds into $\R^{n+1}$ and allows a handle decomposition realizing $P$. To prove that it terminates, it remains to show that such a $W$ exists:
    %
    % We realize the presentation complex $X$ of $P$ as a 2-dimensional finite abstract simplicial complex. It therefore PL-embeds into $\R^{n+1}$ for $n\geq2\cdot2$ by a general position argument \cite[Theorem\~5.4]{rourkesandersonnew}. Hence, there exists a PL-simplicial structure on $\R^{n+1}$ such that $X$ arises as a subcomplex (see \cite[Theorem\~2.14]{rourkesandersonnew}). A regular neighbourhood of $X$ in it has a handle decomposition realizing $P$ by \cite[Proposition\~6.9]{rourkesandersonnew}. Since $\Q^n\subseteq\R^n$ is dense, we may assume that the vertices of this regular neighbourhood have rational coordinates. \comm{TH}{I think this is a bit to quick, I am thinking about it.}
\end{proof}
We will crucially use that $W$ embeds into $\R^{n+1}$ because this restricts how the 2-handles can be attached. In a certain sense, it forces all of them to be \enquote{untwisted}. We use this crucially in the following lemma which is the main idea in Markov's proof\~\cite{Markov} and says that the PL-homeo\-mor\-phism type of a manifold in realizing $P$ detects whether $P$ presents the trivial group. 
\begin{lemma}\label{lem:PtrivialiffS2}
    Let $n\geq4$ and $P$ be a finite group presentation with $l$ relations. Let $W$ be an $(n+1)$-dimensional PL-simplicial manifold that allows a PL-handle decomposition realizing $P$ and PL-embeds into $\R^{n+1}$. Then $P$ presents the trivial group if and only if $\partial W$ is PL-homeomorphic to $\#_l(S^2\times S^{n-2})$.
\end{lemma}
\begin{proof}
    For the \enquote{if}-direction observe that $P$ presents the fundamental group of $W$ and the inclusion $\partial W\to W$ induces an isomorphism on fundamental groups as $W$ has a handle decomposition where all handles have codimension at least 3.

    Now assume that $P$ presents the trivial group. We first prove that $W$ has a PL-handle decomposition consisting of a 0-handle and $l$ many 2-handles. As above, the inclusion $\partial W_2\to W_2$ induces an isomorphism on fundamental groups. By construction, $\pi_1(W_2)$ is presented by $P$, hence $\pi_1(\partial W_2)$ is trivial. This implies that the attaching spheres of the $2$-handles attached in the last bullet point of Definition~\ref {dfn:relalizeP} are homotopic in $\partial W$ to loops that cancel the $1$-handles as in \cite[Lemma 6.4]{rourkesandersonnew}. Since the dimension of $\partial W$ is at least\~4, they are also PL-isotopic to such a loop (compare \cite[Corollary 5.9]{rourkesandersonnew}). It follows that $W$ has a handle decomposition given by attaching the $l$ many 2-handles from the second bullet point of Definition~\ref {dfn:relalizeP} to a $D^{n+1}$.

    Let $A$ be the wedge of spheres formed by the cores of the 2-handles and the cone of their attaching spheres in the 0-handle $D^{n+1}$. After replacing the combinatorial structure on $D^{n+1}$ by the cone of $\partial D^{n+1}$ in the given structure, we may assume that $A\subseteq W$ is a subcomplex. This is still a PL 0-handle as a sequence of stellar subdivisions on stellar welds on $\partial D^{n+1}$ induces one on its cone, hence the PL-homeomorphism type of $W$ remains unchanged. By \cite[Corollary\~3.30]{rourkesandersonnew} $W$ is a regular neighbourhood of $A$ in $\R^{n+1}$ since $W\hookrightarrow\R^{n+1}$ has \mbox{codimension-0}. As a PL-embedding $\bigvee_l S^2\hookrightarrow \R^{n+1}$ is unique up to PL-isotopy by \cite[Theorem\~1]{Husch}, the uniqueness of regular neighbourhoods \cite[Theorem\~3.8]{rourkesandersonnew} yields that $W$ is unique up to PL-homeomorphism. Since the $l$-fold boundary connected sum of $S^2\times D^{n-1}$ is certainly a possibility, the claim follows. 
\end{proof}


%OLD VERSION: In Definition~\ref{def:MP-constr} we associated to each group presentation $P=\langle g_1,\dots,g_k \mid r_1,\dots,r_l\rangle$  a collection $\mathcal{W}^n_P$ of $(n+1)$-dimensional \emph{smooth} manifolds whose fundamental groups are presented by $P$. In the following we will associate to such a presentation $P$ an 
% $(n+1)$-dimensional \emph{PL-simplicial} manifold realizing on element in $\mathcal{W}_P^n$. We do not claim that our construction gives a full algorithm, but we hope that it convinces the reader that one could in principle make this completely algorithmic.
% In the following construction we make use of the following observation: If given two finite abstract simplicial complexes we know for theoretical reasons that a combinatorial isomorphism exists, then 
% one can also find it by going through all (countably many!) possibilities. 



% \begin{construction}\label{construction:xp}
% Let $P=\langle g_1,\dots,g_k \mid r_1,\dots,r_l\rangle$ be a finite presentation of a group. We build an $(n+1)$-dimensional PL-simplicial manifold
% $X_P$ as follows.
% \begin{enumerate}
% \item   We attach $k$ many 1-handles via orientation-preserving simplicial maps 
% to a barycentric subdivision of the simplicial ball
% $\Delta^{n+1}=\{\{0,\dots,n+1\},\mathcal{P}(\{0,\dots,n+1\})\}$. 
% Let $X_1$ be the resulting PL-simplicial manifold. 
% Note that we have a canonical identification $\pi_1(X_1)\cong\langle g_1,\dots,g_k\rangle$.

% \item Possibly after a further barycentric subdivision we can pick disjoint closed oriented simplicial  loops $C_1,\dots,C_l\subset \partial X_1$ 
% such that 
% $C_1,\dots,C_l$ represent the free homotopy classes 
% of $r_1,\dots,r_l\in \pi_1(X_1)\cong \langle g_1,\dots,g_k\rangle$.

% Next we can consider the regular neighbourhoods $N_1,\dots,N_l$
% of $C_1,\dots,C_l\subset \partial X_1$ as in \cite[p.~33]{rourkesanderson} or \cite[Chapter~III, p.~22]{zeeman}. Recall that these can be constructed by taking the double barycentric subdivision of $X_1$ and by considering the union of the stars of all simplices in $N_1,\dots,N_l$.

% It follows from \cite[Theorems~3.8 and 3.11]{rourkesanderson} that each $N_i$ is combinatorially homeomorphic
% to the PL-simplicial manifold $\partial \Delta^2\times \Delta^{n-1}$. It follows from \cite[Lemma~II.3]{Glaser} (applied twice, first to $\Delta^2\times \partial \Delta^{n-1}$ and then to $\Delta^2\times \Delta^n$)
% that for each $i\in \{1,\dots,n\}$ there exists a PL-simplicial manifold $Y_i$ that is 
% combinatorially homeomorphic to $\Delta^2\times \Delta^{n-1}$
% and such that the combinatorial structure 
% on $\partial \Delta^2\times \Delta^{n-1}$ is isomorphic to $N_i$.
% We can now use this isomorphism to attach combinatorial 2-handles to $X_1$ along $N_1,\dots,N_l$. We call the resulting combinatorial  manifold $X_2$.
% \item Finally we  perform the combinatorial boundary connected sum of $X_2$ with $k$ copies of $\partial \Delta^3\times \Delta ^{n-2}$. Here the boundary connected sum is defined as follows:
% We perform two more barycentric subdivisions and then we glue along an $n$-simplex along the boundary.
% \end{enumerate}
% \end{construction}

% \begin{lemma}\label{lem:pl-for-wp}
% Let $P=\langle g_1,\dots,g_k \mid r_1,\dots,r_l\rangle$ be a finite presentation of a group.
% The topological realization of any combinatorial 
% manifold $X_P$ that we constructed in Construction~\ref{construction:xp} is homeomorphic to one of the 
% $(n+1)$-dimensional smooth manifolds that 
% we constructed in Definition~\ref{def:MP-constr}.

% \comm{Marc:}{I think what we really need is the following statement: For any $n\geq4$, there exist an algorithm that\begin{itemize}
%     \item takes as input a group presentation $P$, and
%     \item outputs a PL- simplicial manifold whose underlying PL-manifold is PL-homeomorphic to one element in $\mathcal{W}_P^n$.
% \end{itemize}I think the construction together with the statement of the lemma yield this. }
% \end{lemma}

% \begin{proof}
% The idea of the proof is to follow the logic of 
% Lemma~\ref{lem:welldefined}. The literature that relates smooth constructions to PL-constructions is unfortunately not very extensive. So we will not attempt to provide full details. \comm{Marc:}{Can we maybe add some more details on how this 'should' work?}
% \end{proof} 

% The following lemma is a strengthening of Lemma~\ref{lem:pl-for-wp} and allows us to algorithmically construct PL-simplicial manifolds for all elements in $\mathcal{W}_P^n$. For the proof of Theorem~\ref{thm:main_triangulation_input_and_arbitrary_output} we only need Lemma~\ref{lem:pl-for-wp}, but in the proof of Theorem~\ref{thm:pi_trivial} and Theorem~\ref{thm:unrecognizable_different_categories} we will also use Lemma~\ref{lem:all}.

% \begin{lemma}\label{lem:all}
% There exist an algorithm that   
% \end{lemma}

% \begin{lemma}
% Let $M$ be a compact smooth manifold that is equipped with a smooth simplicial structure $S$ and let $f\colon M\to M$  be a diffeomorphism.
% Then there exist subdivisions $S'$ and $S''$ of $S$ 
% and a simplicial isomorphism $\varphi\colon S'\to S''$
% such that $|\varphi|$ is homeotopic to $f$.    
% \end{lemma}

% \begin{proof}
% This lemma follows from the uniqueness of smooth simplicial structures on smooth manifolds (see \cite[Theorem~10.5]{Munkres1966}) applied to the smooth simplicial structures $S$ and $f^*S$. 
% \end{proof}

% In our application we will apply the lemma to $M=D^n\times S^1$, 
% a fixed smooth simplicial structure on $D^n\times S^1$ and the diffeomorphism $D^n\to D^n\times S^1$ that is given by twisting along $S^1$. 


%==========================================
\section{Unsolvability of the decision problem}
In this section, we will give the proof of Theorem\~\ref{thm:main_triangulation_input_and_arbitrary_output}. For that we first recall that the triviality problem for finitely presented groups is unsolvable.

\begin{theorem}[Novikov\~\cite{Novikov-algo-insolv}]\label{thm:triviality_problem}
    There exists no algorithm that
    \begin{itemize}
        \item takes as input a finite group presentation, and
        \item outputs whether or not it presents the trivial group.\qed
    \end{itemize}
\end{theorem}


\begin{proof}[Proof of Theorem\~\ref{thm:main_triangulation_input_and_arbitrary_output}]
We consider an integer $n\geq 4$ and one of the following equivalence relations on $n$-dimensional PL-manifolds:
\begin{enumerate}
    \item being PL-homeomorphic,
    \item being homeomorphic,
    \item being homotopy equivalent,
    \item admitting isomorphisms between all homotopy groups,
    \item admitting an isomorphism between the fundamental groups.
\end{enumerate}
We will argue by contradiction. We assume that there exists an algorithm that
\begin{itemize}
    \item takes as input $n$-dimensional PL-simplicial manifolds $M_1$ and $M_2$ %that admit a smooth structure, and 
    \item outputs whether or not $|M_1|$ is equivalent to $|M_2|$.
\end{itemize}
We claim that the existence of such an algorithm would imply the existence of an algorithm deciding if a group is trivial and thus contradict Theorem\~\ref{thm:triviality_problem}. 

Indeed, let $P$ be a finite group presentation. Using Lemma\~\ref{lem:algrelP}, we can algorithmically construct a PL-simplicial manifold $W$ which PL-embeds into $\R^{n+1}$ and allows a handle decomposition realizing $P$. Now create a PL-simplicial manifold $T$ realizing $\#_l(S^2 \times S^{n-2})$ where $l$ is the number of relations in $P$ and use the potential algorithm to decide if $\partial W$ and $T$ represent equivalent spaces. By Lemma\~\ref{lem:PtrivialiffS2} this is the case if and only if $P$ presents the trivial group.
% Indeed, let $P$ be a finite presentation of a group $G$. Using Lemma\~\ref{lem:pl-for-wp} we can algorithmically construct a PL-simplicial manifold $T$ realizing $\partial W$ where $W$ is one of the manifolds in $\mathcal{W}_P^n$, see Definition\~\ref{def:construction}. Let $l$ be the number of relations in $P$. By Lemma\~\ref{lem:P}, an integer $m\geq0$ such that $\partial W$ is diffeomorphic (and hence equivalent) to 
% $$\#_{l-m} (S^2 \times S^{n-2}) \#_m (S^2 \ttimes S^{n-2})$$
% exists if and only if $G$ is isomorphic to the trivial group. 
%
% Next, we create PL-simplicial manifolds $T_0,\ldots,T_l$ realizing all the orientable smooth manifolds $\#_{l-m} (S^2 \times S^{n-2}) \#_m (S^2 \ttimes S^{n-2})$ for $m=0,\ldots,l$. For each $(T_i)_{i=0,\ldots,l}$, we use our potential algorithm to decide if $T$ and $T_i$ represent equivalent spaces. If this is the case for one $T_i$, then $G$ is isomorphic to the trivial group, and if $T_i$ and $T$ represent non-equivalent spaces for all $i=0,\ldots,l$, then $G$ is not isomorphic to the trivial group by Lemma\~\ref{lem:P}.
\end{proof}

\section{Unrecognizable manifolds}\label{sec:unrecognizable}

%In this section we prove Theorem\~\ref{thm:pi_trivial} and\~\ref{thm:unrecognizable_different_categories}. For that we will need the following strengthening of Theorem\~\ref{thm:triviality_problem}.
The proof of Theorems\~\ref{thm:pi_trivial} and\~\ref{thm:unrecognizable_different_categories} is obtained in the same way from the following strengthening of Theorem\~\ref{thm:triviality_problem}.

\begin{theorem}[Gordon {\cite[Theorem 2.2]{Gordon}}] \label{thm:Adian_Rabin_set}
For every $l\geq13$, there exists an \emph{Adian--Rabin set} $S$ with $l$ relations, i.e.\ a set of finite group presentations with exactly $l$ relations, such that there exist no algorithm that
\begin{itemize}
    \item takes as input a presentation $P$ from $S$, and
    \item outputs whether or not $P$ presents the trivial group.\qed
\end{itemize}
\end{theorem}

% \begin{proof} [Proof of Theorem\~\ref{thm:pi_trivial} and\~\ref{thm:unrecognizable_different_categories}]
% By Theorem\~\ref{thm:Adian_Rabin_set} there exist an\~$l\in \N$ such that there exists an Adian--Rabin set $S$ with\~$l$ relations. Choose a natural number $n\geq4$ and any of the equivalence relations used in Theorem\~\ref{thm:pi_trivial} or Theorem\~\ref{thm:unrecognizable_different_categories}. Similar as in the proof of Theorem\~\ref{thm:main_triangulation_input_and_arbitrary_output} we argue by contradiction and assume that there exist an algorithm that
% \begin{itemize}
%     \item takes as input a PL-simplicial $n$-dimensional manifold $M$ %that admits a smooth structure, and 
%     \item outputs whether or not $|M|$ is equivalent to $\#_l (S^2\times S^{n-2})$.
% \end{itemize}
% From this we will construct an algorithm that takes as input presentations from $S$ and outputs if it represents the trivial group, contradicting Theorem\~\ref{thm:Adian_Rabin_set}.

% For that let $P$ be a finite group presentation in $S$. 
% From Lemma\~\ref{lem:P_strong} we know that at least one manifold in $\mathcal{W}_P^n$ is diffeomorphic (and hence equivalent) to $\#_l (S^2\times S^{n-2})$ if and only if $P$ presents the trivial group.
% Using Lemma\~\ref{lem:all} we can algorithmically construct PL-simplicial manifolds $M_1,\ldots,M_{k}$ representing all elements in $W_P$. Similarly we can also build a PL-simplicial manifold $M$ representing $\#_l (S^2\times S^{n-2})$. For each $(M_i)_{i=1,\ldots,k}$, we use our potential algorithm to decide if $|M|$ and $|M_i|$ represent equivalent spaces. If this is the case for one $M_i$, then $P$ presents the trivial group, and if $M_i$ and $M$ represent non-equivalent spaces for all $i=1,\ldots,k$, then $P$ does not represent the trivial group by Lemma\~\ref{lem:P_strong}.
% \comm{Marc}{I think Stefan's proof of Lemma\~\ref{lem:all} will actually produce $2^l$ instead of $k$ manifolds. I.e. it will overcount in some situations. But let's maybe see what he writes there. But I think we have to go back to replace $k$ here by $2^l$.}
% \end{proof}


%============================================







\let\MRhref\undefined
\bibliographystyle{hamsalpha}
\bibliography{bib.bib}

\end{document}


\section{old}

% We reduce the above diffeomorphism problem to the triviality problem for groups using the following construction.

% \begin{definition}
% \label{def:MP-constr}
% Let $P=\langle g_1,\dots,g_k \mid r_1,\dots,r_l\rangle$ be a finite group presentation of a group. Starting with $P$, we build a set $W_P$ of at most $2^l$ compact smooth $(n+1)$-manifolds, where each manifold arises from the following construction.
% \begin{enumerate}
% \item \label{def:MP-constr-2}  We attach $k$ many 1-handles via orientation-preserving attaching maps to $D^{n+1}$. We call the resulting smooth manifold $W_1$. Note that there exists a canonical isomorphism $\pi_1(W_1)=\langle g_1,\dots,g_k\rangle$. The inclusion induced map $\pi_1(\partial W_1)\to \pi_1(W_1)$ is an isomorphism because we only attach handles of codimension at least 
% three.
% \item \label{def:MP-constr-3} We attach $l$ many 2-handles to $W_1$ along attaching maps $\varphi_1,\dots,\varphi_l\colon\partial D^2\times D^{n-2}\to \partial W_1$ such that for each $i\in \{1,\dots,l\}$ the curve $\varphi_i(\partial D^2\times\{0\})$  represents $r_i\in \langle g_1,\dots,g_k\rangle$ under the above isomorphism. In dimension $n\geq4$ homotopy of curves implies smooth isotopy. It follows from this observation and the uniqueness of tubular neighbourhoods  (see \cite{friedlmonster}) \comm{TH}{more precise referece, also newer version of notes} that for this attachment, there are (up to smooth isotopy) two possible choices for the 
% framing of a single $2$-handle. \comm{MK}{This is because of $\pi_1(SO(3))=\Z_2$? Maybe this could be mentioned.} We build a manifold for each such choice. We call the resulting smooth manifold $W_2$.
% \item \label{def:MP-constr-4} Finally we attach $k$  
% additional 2-handles to $D^{n+1} \subseteq W_2$ such that the attaching curves are null-homotopic with the trivial framing. Note that in this standard situation of $D^{n+1}$, it is possible to speak of the trivial framing by demanding that $D^5\cup h_2$ is a product.
% \end{enumerate}
% Note that because of the choices in the third step, we end up with a set of at most $2^l$ compact smooth $(n+1)$-manifolds.
% \end{definition}

% \begin{lemma}
% Let $P$ be a finite group presentation of the trivial group and $W\in W_P$. There exists $k\in\N_0$ such that $\partial W$ is diffeomorphic to 
%     \[
%     \#_{l-k} (S^2 \times S^{n-2})\#_k(S^2\ttimes S^{n-2})
%     \]
%     where $l$ is the number of relators in $P$.
% \end{lemma}
% \begin{proof}
%     In Step \ref{def:MP-constr-2} of \cref{def:MP-constr} \comm{TH}{I dont knot why this is quoted as a proposition} we already pointed out that $\pi_1(W)\cong\pi_1(\partial W)$. By construction, $P$ presents $\pi_1(\partial W)\cong\pi_1(W)$ which are therefore trivial. 

%     Consider the trivially framed 2-handles attached in Step \ref{def:MP-constr-4} in the construction of $W$. 
%     Because $\pi_1(\partial W_2)\cong\pi_1(\partial W)$ is trivial, the attaching sphere of such a $2$-handle is homotopic, hence isotopic, to a loop that cancels one of the $1$-handles 
%     attached in Step \ref{def:MP-constr-2}. 
%     Thus, $W$ is diffeomorphic to $D^{n+1}$ with $l$ 
%     many $2$-handles attached. In Step \ref{def:MP-constr-3} we discussed that there are precisely two choices of framing for each 2-handle, with in the boundary lead to a summand of $S^2 \times S^{n-2}$ or $S^2\ttimes S^{n-2}$.
% \end{proof}

\begin{lemma}
\label{lem:MP-trivtangentbundle}
Among the $2^l$ choices in the above definition, there is at least one choice that leads to a smooth manifold with trivial tangent bundle.
\end{lemma}

\begin{proof}
It suffices to show that in the above construction at least one choice leads to a codimension-zero submanifold of  $\R^5$.
We attach the 1-handles in the standard way to obtain a codimension-zero submanifold $W_1$ of $\R^5$. Let $X_1$ be  the complement of the interior of $W_1$ viewed as a submanifold of $S^5=\R^5\cup \{\infty\}$. Note that $X_1$ admits a handle decomposition with one 0-handle and $k$ 3-handles. In particular, we see that $X_1$, and thus also $X_1\cap \R^5$ is simply connected.

\comm{MK}{Maybe we can make the above a bit more precise. Something like: For every $1$-handle $h^i_1$ of $W_1$ we attach a $2$-handle $h_2^i$ to $W_1$ such that the attaching sphere of the $2$-handle $h_2^i$ intersects the belt sphere of the $1$-handle $h_1^i$ transversely in a single point. Thus the handles cancel and we get a handle decomposition of $D^5$ to which we can glue a $5$-handle to obtain $S^5$. This describes an embedding of $W_1$ into $S^5$. By considering the dual handle decomposition, the exterior of $W_1$ in $S^5$ can be obtained by attaching $k$ $3$-handles to a $0$-handle and thus is simply connected. By seeing $S^5$ as $\R^5\cup \{\infty\}$, we also get an embedding of $W_1$ into $\R^5$ with simply connected exterior $X_1$.}

Now let $C_1,\dots,C_l$ in $\partial W_1$ be smoothly embedded curves that correspond to $r_1,\dots,r_l$. Since 
$X_1\cap \R^5$ is simply connected and since $\dim(X_1)=5>2\cdot 2$
we see that there exist proper smooth embeddings $\varphi_1,\dots,\varphi_l\colon D^2 \to X_1\cap \R^5 $
such that for each $i$ we have $\varphi_i(S^1)=C_i$.
For the proper submanifolds $\varphi_i(D^2)\subseteq X_1$, $i=1,\dots,l$ we pick tubular neighbourhoods. Since $D^2$ and $\R^5$ are orientable, we see that the tubular neighbourhoods are trivial. We now use these trivial tubular neighbourhoods to attach the $l$ 2-handles
and to form $W_2$, which \comm{MK}{corresponds to one of the $2^l$ choices and} is now by construction a codimension-zero submanifold of~$\R^5$.

Finally the $k$ remaining 2-handles are attached in the standard \comm{MK}{I am not sure what standard way means. Maybe we better just say that the $2$-handle attachment from point (4) can be done in $\R^5$.} way and this can evidently be done in $\R^5$.
\end{proof}

\begin{lemma}
\label{lem:Ptrivial-iffMP-S2}
Let $P=\langle g_1,\dots,g_k\,|\, r_1,\dots,r_l\rangle$ be a finite group presentation \comm{MK}{of a group $G$}. 
The group~$P$ \comm{MK}{$G$ or $P$ represents the trivial group. Similar below.} is the trivial group if and only if
at least one of the manifolds in $M_P$ \comm{MK}{$M_P$ is the collection of the $2^l$ manifolds constructed above?}  has boundary diffeomorphic to $\sumStwo{l}$. 
\end{lemma}

\begin{proof}
By construction, we have~$\pi_1(M) \cong P$
for every~$M\in M_P$. Moreover, because
we only attach handles of codimension at least 
three, we have $\pi_1(\partial M) \cong \pi_1(M)$.
Hence, the `only if' part of the statement is clear. \comm{MK}{We could refer here to point (2) of the definition.}

Now, suppose that~$P$ is the trivial group.
By Lemma~\ref{lem:MP-trivtangentbundle}, 
we can find~$M \in M_P$ with trivial tangent bundle. 

Consider Step~\ref{def:MP-constr-4} in the construction (see Definition~\ref{def:MP-constr}). 
Because~$\pi_1(W_2) \cong P \cong 1$, the attaching
sphere of a $2$-handle is homotopic, hence isotopic,
to a loop that cancels one of the $1$-handles 
attached in Step~\ref{def:MP-constr-2}. 
Thus, $M$ is diffeomorphic to~$D^5$ with~$l$ 
many $2$-handles attached. Hence, it is
diffeomorphic to $\mathop\natural_m S^2 \times D^3 \natural \mathop\natural_{l-m} S^2 \widetilde\times D^3$
for some~$m\in \{0,\dots,l\}$.  
%\comm{MU}{Should we add more details why? MP: could say, $M$ has trivial tangent bundle hence $\partial M$ does, but there are no $S^2 \wt{\times} S^2 = \CP^2 \# \ol{\CP}^2$ summands because we know that $\CP^2$ has odd intersection form so cannot be spin, so $w_2 \neq 0$.} 
%\comm{MU}{The triviality of the 
%tangent bundle does not pass to
%the boundary (Consider \eg a sphere),
%but stable triviality does.  MP: That's what I meant to say, and also what you need.}
Because~$M$ has trivial tangent bundle, $\partial M$ has stably trivial tangent bundle. Hence there are no $S^2 \wt{\times} S^2 = \CP^2 \# \ol{\CP}^2$ summands, because if there were, there would be a $\CP^2$ summand. But $\CP^2$ has odd intersection form, so is not spin, and hence  $w_2(S^2 \wt{\times} S^2) \neq 0$, contradicting that the tangent bundle of $\partial M$ is stably trivial. 
Whence~$m=l$, and thus $\partial M$ is diffeomorphic to~$\sumStwo{l}$.
\comm{MK}{Here is a slight reformulation (and in particular not using a contradiction): It follows that the boundary $\partial M$ is diffeomorphic to $\#_m S^2 \times S ^2 \#_{l-m} S^2 \widetilde\times S^2$. Since $\partial M$ embedds into $\R^5$ it has stably trivial normal bundle and thus its second Stiefel--Whitney class vanishes $w_2(\partial M)=0$. On the other hand, it follows that $S^2 \widetilde\times S^2$ has odd intersection form, so it is not spin and hence $w_2(S^2 \widetilde\times S^2)\neq0$. Whence~$m=l$, and thus $\partial M$ is diffeomorphic to~$\sumStwo{l}$.
}
\end{proof}

\begin{proof}[Proof of Theorem~\ref{thm:unsolv-diffeo}]
    Fix~$l\in \N$ such that there exists an Adian--Rabin set with~$l$ relations, i.e.\
    the triviality problem for group presentations with exactly~$l$ relations is 
    unsolvable. \comm{MU}{needs ref that such a thing exists MK: Maybe we should also properly define what this means.} Assuming the existence
    of an algorithm~$A$ as in the theorem, we show solvability of the triviality problem.
\comm{MT}{If you want an explicit reference, you can use Theorem 2.2 in~\cite{Gordon}. It provides an Adian-Rabin set with $l=13$. If you want to optimise to get the best possible $l$, then I think that it is possible to get $l=11$ but it requires a bit more work/explanation: In my paper, I construct an Adian-Rabin set of deficiency $9$. Because it is constructed via Gordon-Miller approach, this set can be transformed into an Adian-Rabin set where each presentation has 11 relations and 2 generators. (But I was not doing so because I do not need it.) This is mentioned in Miller's paper ``Decision problems for groups—survey and reflections'' and if I remember well also reproduced in~\cite{Chernavsky}. Of course, the question is whether it is worth optimising.}

    Given a group presentation~$P=\langle g_1,\dots,g_k\,|\, r_1,\dots,r_l\rangle$ with
    $l$ relations, construct the set of 5-manifolds~$M_P$.
    For each $M\in M_P$, query~$A$ whether $\partial M$ is diffeomorphic to 
    $\sumStwo{l}$. \comm{MK}{maybe here we should mention or explain that we get a Kirby diagram of $\partial M$ from the data of $P$ so that we can apply the algorithm. We could mention that the same argument works for any other presentation of $\partial M$ that can be extracted from the above construction.} By Lemma~\ref{lem:Ptrivial-iffMP-S2}, if the answer
    is affirmative in one case, this certifies that~$P$ presents the trivial group.
    Otherwise, if the answer is always negative, $P$ does not present the trivial group. \comm{MK}{Maybe here we should more clearly explain why this is a contradiction. In particular we should explain the use of the Adian--Rabin set (and why this is needed) and cite the result why this is a contradiction.} \comm{SD}{I second MK}
\end{proof}

\comm{MK}{In fact, I am confused now. It seems to me that we only show that the diffeomorphism problem is not solvable. In principle it could be that there exist an algorithm deciding homeomorphism of smooth manifolds. Is that correct? }
\comm{SF}{The argument does both. We should mention this.}

